\section{Background}

\subsection{Application Boundary and Evidence Scale}

A quantum application wraps a classical input in a circuit and returns samples, observables, kernel entries, or energies to a classical workflow. The kernel is only one stage. Encoding, repeated execution, measurement aggregation, optimizer/search iterations, and postprocessing remain part of application time and quality.

The four families exercise different paths. A quantum neural network or variational classifier (QNN/VQC) trains a parameterized feature circuit through classical feedback; QKernel repeatedly evaluates feature-map circuits to construct a kernel matrix; QFeature executes a fixed circuit and trains a classical head. Chem. uses a VQE ansatz and grouped Hamiltonian measurements. Opt. uses QAOA phase/mixer circuits and sampled objectives. Sim. maps Hamiltonian terms to Trotterized evolution. Their native counterparts are classifiers, eigensolvers, exact or heuristic optimization, and Krylov time evolution. Each pair shares an input and reports quality in accuracy, Hartree, objective ratio, or observable error.

A case crosses the application boundary only when
\[
T_{\mathrm{qhw}}<T_{\mathrm{native}}
\quad\text{and}\quad
\Delta Q\leq\epsilon_w,
\]
where $\epsilon_w$ is the workload's natural-unit tolerance. More shots may narrow estimator error, while a deeper ansatz, search, or Trotter product may reduce algorithmic error; both increase hardware demand. Quality recovery is therefore measured jointly with cost rather than represented by a free scalar.

Evidence scale and circuit width are distinct. Controlled 4--20-qubit records permit exact native answers and natural-unit quality closure. Distributed 36--40-qubit state-vector runs establish HPC simulation capacity only. Deployment-facing native workloads tighten $T_{\mathrm{native}}$ but do not inherit the controlled circuit's quality. A projected point becomes an application-level hardware target only when finite-shot output, the complete application loop, dependency schedule, reliability, and physical model attach to the same record.

\subsection{Reliability-Constrained Critical Path}

Surface-code execution couples reliability, time, and space. We use the QDK-aligned logical-error envelope
\[
p_L(d)=0.03\left(p/0.01\right)^{\lfloor(d+1)/2\rfloor},
\]
and choose the smallest odd distance whose logical-volume union bound fits the allocated application failure budget~\cite{azurequantumresourceestimatordocs}. A strict contract requires the complete estimator's logical, synthesis, and distillation failures to fit 1\% or 0.1\%. A separately reported estimator-tolerant contract permits corrupted shots only when direct finite-shot output provides enough measured headroom; it does not assume that every shot must be error-free.

Arbitrary rotations add a second reliability constraint. The leading Ross--Selinger synthesis cost is $3\log_2(1/\epsilon_{\mathrm{rot}})$ T states, where the per-rotation precision budget follows from the number of rotations in the application loop~\cite{rossselinger2016}. A fixed 16- or 32-T assumption is invalid whenever it exceeds that budget. Factory output is then priced with LSQCA's 176-cell surface-code factory producing one state every 15 code beats, while QDK's Litinski protocols supply the corresponding physical-qubit footprint~\cite{lsqca2025,azurequantumresourceestimatordocs}. These are separate throughput and space envelopes; they are not multiplied together as independent speedups.

For logical work $T_{\mathrm{log}}$, decoding/reaction $T_{\mathrm{dec}}$, and factory supply $T_{\mathrm{fac}}$, unavailable contention is bracketed by
\[
T_{\mathrm{core}}(\rho)=\max_i T_i+
\rho\left(\sum_i T_i-\max_i T_i\right),
\quad \rho\in[0,1].
\]
Decoder service and the communication/control path remain separate because reaction time can require correction storage and additional protected memory~\cite{khalid2025reaction}. Host-visible feedback is serial only when the recorded loop requires it. Likewise, useful shot lanes are capped by dependency-ready repetitions, not physical lane count alone. A component design such as LSQCA, BOSS, S-SYNC, Pinball, or Cyclone can replace only matched core-cell, movement, decoder, or topology events~\cite{lsqca2025,boss2025,ssync2025,pinball2026,cyclone2026}. This replacement rule is what lets an application boundary identify when a published mechanism is first-order and when it is not.
